\begin{table}[t]
  \centering
  \caption{Same-dataset comparison on LECSEG-30. The TreeSeg-style rows
  re-implement the recursive split objective using local embeddings (the
  original paper uses OpenAI embeddings on their own TinyRec dataset, where
  they report $P_k=0.367$; that number is \emph{not} directly comparable here
  because the datasets differ). Lower is better for $P_k$/WD; higher for F1@2.}
  \label{tab:same_dataset_comparison}
  \small
  \resizebox{\linewidth}{!}{%
  \begin{tabular}{lrrrl}
    \toprule
    Method & $P_k\downarrow$ & WD$\downarrow$ & F1@2$\uparrow$ & Interpretation \\
    \midrule
    Balanced LOO selector & 0.3588 & 0.3739 & 0.0893 & Official LECSEG best \\
    Cross-model conservative & 0.3713 & 0.3764 & 0.0237 & Best single global method \\
    TreeSeg-style (MPNet, LECSEG-30) & 0.4320 & 0.4673 & 0.1733 & Higher F1, worse $P_k$/WD \\
    TreeSeg-style (E5-large, LECSEG-30) & 0.4322 & 0.4654 & 0.1576 & Same-dataset comparator \\
    TreeSeg-style (BGE-large, LECSEG-30) & 0.4399 & 0.4780 & 0.1643 & Same-dataset comparator \\
    \midrule
    TreeSeg paper (on TinyRec)\textsuperscript{$\dagger$} & 0.367 & --- & --- & Different dataset; indicative only \\
    \bottomrule
  \end{tabular}%
  }
  \vspace{0.3em}
  \small\textsuperscript{$\dagger$}~Gklezakos et al.\ (2024), arxiv:2407.12028. Evaluated on 21 self-recorded lectures (TinyRec), not LECSEG-30.
\end{table}
